\documentclass[10pt,a4paper]{amsart}
\usepackage[margin=19mm]{geometry}
\usepackage{amsmath,amssymb,mathtools}
\usepackage[scr=rsfs,cal=euler]{mathalfa}
\usepackage{xfrac}
\usepackage[unicode,hidelinks,bookmarks=false]{hyperref}
\newcommand{\ie}{i.e.\ }
\newcommand{\cf}{cf.\ }
\newcommand{\resp}{respectively\ }
\newcommand{\Subsection}{\subsection}
\providecommand{\nobreakdash}{\nobreak\hbox{-}}
\setlength{\emergencystretch}{2em}
\multlinegap0pt
\usepackage{amssymb}
\usepackage{mathtools}
\newtheorem{theorem}{Theorem}
\newtheorem{lemma}{Lemma}
\newtheorem{corollary}{Corollary}
\newcommand{\Del}{\Delta}
\newcommand{\DelTwo}{\Delta_2}
\newcommand{\cH}{\mathcal H}
\newcommand{\cJ}{\mathcal J}
\newcommand{\cQ}{\mathcal Q}
\newcommand{\cQlin}{\mathcal Q_{\mathrm{lin}}}
\newcommand{\ta}{\tau}
\title{An Improved Lower Bound for the Erd\H{o}s--Lov\'asz Cover Number Problem}
\author{Varun Sivashankar}
\date{Source: arXiv:2606.24878v2 (CC BY 4.0)}
\begin{document}
\maketitle

\noindent\emph{Source and licence.} An Improved Lower Bound for the Erd\H{o}s--Lov\'asz Cover Number Problem. Version arXiv:2606.24878v2 (CC BY 4.0), distributed by arXiv under the Creative Commons Attribution 4.0 International licence, http://creativecommons.org/licenses/by/4.0/, as recorded in the arXiv licence metadata for this exact version and shown on the abstract page https://arxiv.org/abs/2606.24878v2. Full licence notice: \texttt{licenses/arxiv-2606.24878-CC-BY-4.0.txt}.\par\smallskip

\noindent\textit{Editorial scope.} Counted unit: the article's theorem thm:main (source0.tex lines 87--100). The article's complete original proof of it is delivered with it (source0.tex lines 502--697, one \texttt{proof} environment of 196 lines, which the article places away from the statement). No mathematics is written, repaired or re-derived: the statement and the proof are the article's own lines, in the article's own notation, and no retained line is substituted or re-flowed.  Retained uncounted as the source-local context the proof needs: lines 219--225; lines 473--484.  Nothing was omitted from any retained span, so the package carries no omission record.  No retained line contains a citation, so the wrapper carries no bibliography and no bibliography entry.  No retained line contains a figure environment, an image-inclusion command or drawing code, so no image is delivered and the package declares no asset dependency.  Editorial formatting only, in addition: an AMS \texttt{amsart} preamble that restates the article's own declarations and macros used by the retained lines, namely the environments \texttt{theorem}, \texttt{lemma}, \texttt{corollary} on separate counters, together with the standard packages those lines need.  The retained environments print in this document as Theorem 1, Lemma 1, Corollary 1 in order of appearance, whereas the article prints the same items with the numbering of its own full document.  The complete native source of the article as distributed in the arXiv e-print for this exact version is bundled unchanged as \texttt{sources/203561/source0.tex}.
\begin{lemma}\label{lem:residual-cover}
Let $\cJ$ be an $r$-uniform intersecting hypergraph with $q$ edges and
$\Del(\cJ)\le 3$.  Then
\[
  4\ta(\cJ)\le q+r+4.
\]
\end{lemma}

\begin{corollary}\label{cor:kahn-rank-four}
For every $\eta>0$, there is $D_0$ such that the following holds.  Let $G$ be
a $4$-uniform hypergraph with $\Del(G)=D\ge D_0$ and
$\DelTwo(G)\le 1$.  Then
\[
  \chi'(G)\le (1+\eta)D.
\]
Consequently,
\[
  \nu(G)\ge \frac{|E(G)|}{(1+\eta)D}.
\]
\end{corollary}

\begin{theorem}\label{thm:main}
The following bounds hold.
\begin{enumerate}
\item[(i)] For every positive integer $r$,
\[
  g(r)\ge 3r-4.
\]
\item[(ii)] For every $\varepsilon>0$, there is $r_0$ such that, for every
$r\ge r_0$,
\[
  g(r)\ge \left(\frac{41-\sqrt{19}}{12}-\varepsilon\right)r.
\]
\end{enumerate}
\end{theorem}

\begin{proof}[Proof of Theorem~\ref{thm:main}(ii)]
Put
\[
  \beta_0:=\frac{4+\sqrt{19}}{3},
  \qquad
  c:=\frac{41-\sqrt{19}}{12}.
\]
Fix $\varepsilon>0$, and choose $\eta$ so that
\[
  0<\eta\le
  \min\left\{1,\left(\frac{\varepsilon}{2c}\right)^2\right\}.
\]
We prove that $m\ge (c-\varepsilon)r$ for all sufficiently large $r$.

Let $\cH$ be an $r$-uniform intersecting hypergraph with $\ta(\cH)=r$ and
$m=|E(\cH)|$.  Repeating the peeling argument with threshold $5$, peel
vertices of current degree at least $5$.  Let $S$ be the set of peeled
vertices, let $k:=|S|$, and let $\cJ$ be the remaining hypergraph.  Put
$q:=|E(\cJ)|$.  Then $\Del(\cJ)\le 4$, and the same covering argument gives
\begin{equation}\label{eq:five-peeling}
  m\ge q+5\bigl(r-\ta(\cJ)\bigr).
\end{equation}
Also, if $q>0$, then fixing one edge of $\cJ$ and using $\Del(\cJ)\le 4$
gives
\begin{equation}\label{eq:q-upper-four}
  q\le 3r+1.
\end{equation}
The same inequality is trivial when $q=0$.

We first prove a cover bound that reuses Lemma~\ref{lem:residual-cover}.  Let
$\cQ$ be the auxiliary $4$-uniform hypergraph on vertex set $E(\cJ)$ whose
edge-copies are the four-element blocks $B_v$ coming from vertices $v$ of
$\cJ$ of degree $4$.  Let $M$ be a maximal matching in $\cQ$, say
$|M|=a$, and let $U$ be the set of edges of $\cJ$ not covered by the matched
auxiliary edges.  The $a$ corresponding vertices of $\cJ$ cover the
$4a$ edges outside $U$.  Moreover, the subhypergraph of $\cJ$ induced by $U$
has maximum degree at most $3$, since a degree-four vertex inside $U$ would
give an auxiliary edge disjoint from $M$.  Lemma~\ref{lem:residual-cover}
therefore gives
\[
  \ta(\cJ)\le a+\frac{q-4a+r+4}{4}
  =\frac{q+r+4}{4}.
\]
Substituting this into~\eqref{eq:five-peeling} gives
\begin{equation}\label{eq:first-kahn-bound}
  m\ge \frac{15r}{4}-\frac q4-5.
\end{equation}
Since $c=15/4-\beta_0/4$, if $q\le\beta_0r$, then
~\eqref{eq:first-kahn-bound} gives $m\ge cr-5$, which is at least
$(c-\varepsilon)r$ for all sufficiently large $r$.  We may therefore assume
$q>\beta_0r$.  In particular, $q\to\infty$ and $q=\Theta(r)$,
by~\eqref{eq:q-upper-four}.

It remains to prove a second bound, useful when $q$ is large.  Let $x_i$ be
the number of vertices of degree exactly $i$ in $\cJ$, for $1\le i\le 4$.
Then
\begin{equation}\label{eq:degree-four-incidences}
  x_1+2x_2+3x_3+4x_4=qr.
\end{equation}
For distinct $A,B\in E(\cJ)$, let $\lambda_{AB}$ be the number of
degree-four vertices of $\cJ$ lying in $A\cap B$, and define
\[
  X:=\sum_{\{A,B\}\subseteq E(\cJ)}(\lambda_{AB}-1)_+.
\]
We claim that $\cQ$ has a $4$-uniform subhypergraph $\cQlin$ with
$\DelTwo(\cQlin)\le 1$ and
\begin{equation}\label{eq:rank-four-linear-count}
  |E(\cQlin)|\ge \binom q2-\frac{5qr}{4}.
\end{equation}
Indeed, let $z$ be the number of pairs $\{A,B\}$ with $\lambda_{AB}=0$.
For each such pair choose a vertex in $A\cap B$.  This chosen vertex has
degree $2$ or $3$: it cannot have degree $1$, since it lies in both $A$ and
$B$, and it cannot have degree $4$, by the definition of $z$.  A degree-two
vertex lies in one pair of edges of $\cJ$, while a degree-three vertex lies in
three pairs.  Thus
\[
  z\le x_2+3x_3.
\]
Every degree-four vertex of $\cJ$ lies in six pairs of edges, and hence
contributes $6$ to the sum of the pair-codegrees.  Since $\lambda_{AB}>0$
for exactly $\binom q2-z$ pairs,
\[
  X
  =6x_4-\left(\binom q2-z\right)
  \le 6x_4-\binom q2+x_2+3x_3.
\]
Hence
\[
  x_4-X\ge \binom q2-x_2-3x_3-5x_4
  \ge \binom q2-\frac{5qr}{4},
\]
where the last inequality follows from~\eqref{eq:degree-four-incidences}.
Starting with $\cQ$, whenever a pair $\{A,B\}$ of auxiliary vertices has
current codegree $d\ge2$, delete one entire auxiliary edge-copy $B_v$
containing $\{A,B\}$.  The chosen pair's excess drops from $d-1$ to $d-2$,
and no other pair-excess can increase.  Thus the total excess drops by at
least $1$ per deletion.  Initially it is $X$, so at most $X$ edge-copies are
deleted.  At termination, the remaining subhypergraph $\cQlin$ has
$\DelTwo(\cQlin)\le1$ and at least $x_4-X$ edge-copies.  This
proves~\eqref{eq:rank-four-linear-count}.

Write
\[
  e:=|E(\cQlin)|,\qquad D:=\Del(\cQlin).
\]
Since $q>\beta_0r$ and $\beta_0>5/2$,
~\eqref{eq:rank-four-linear-count} gives $e=\Omega(r^2)$.
Consequently, $D\ge 4e/q=\Omega(r)$.  For all sufficiently large $r$,
Corollary~\ref{cor:kahn-rank-four} therefore yields a matching in $\cQlin$
of size
\[
  a\ge \frac{e}{(1+\eta)D}.
\]
The corresponding $a$ vertices of $\cJ$ cover $4a$ distinct edges; pairing
the remaining edges gives
\[
  \ta(\cJ)
  \le a+\left\lceil\frac{q-4a}{2}\right\rceil
  \le \frac q2-\frac{e}{(1+\eta)D}+1.
\]

The above bound is strong when $D$ is small.  If $D$ is large, we can bound
the cover number in a different way.  Choose a vertex of $\cQlin$ of degree
$D$.  Since $\DelTwo(\cQlin)\le1$, the $D$ auxiliary edges through it have no
other vertices in common, and hence their union has size $1+3D$.  The
corresponding $D$ vertices of $\cJ$ cover those $1+3D$ edges.  Pairing the
remaining edges gives
\[
  \ta(\cJ)
  \le D+\left\lceil\frac{q-1-3D}{2}\right\rceil
  \le \frac q2-\frac D2.
\]
Combining the two covers and using the geometric mean,
\begin{equation}\label{eq:matching-star-cover}
  \ta(\cJ)
  \le \frac q2-
  \max\left\{\frac D2,\frac{e}{(1+\eta)D}\right\}+1
  \le \frac q2-\sqrt{\frac{e}{2(1+\eta)}}+1.
\end{equation}
Substituting~\eqref{eq:rank-four-linear-count} into
~\eqref{eq:matching-star-cover}, and then using~\eqref{eq:five-peeling},
gives
\[
  m\ge
  5r-\frac{3q}{2}
  +5\sqrt{\frac{1}{2(1+\eta)}
  \left(\binom q2-\frac{5qr}{4}\right)}-5.
\]

To finish, write $\beta=q/r$ and define
\[
  F(\beta):=
  5-\frac{3\beta}{2}
  +5\sqrt{\frac{2\beta^2-5\beta}{8}}.
\]
For all sufficiently large $r$, we have
$\beta/(4r)\le(2\beta^2-5\beta)/8$, since $\beta>\beta_0>5/2$.
Using $\sqrt{x-y}\ge\sqrt{x}-\sqrt y$ for $x\ge y\ge0$, first with
$x=(2\beta^2-5\beta)/8$ and $y=\beta/(4r)$, and then in
\[
  \frac1{\sqrt{1+\eta}}
  =\sqrt{1-\frac{\eta}{1+\eta}}\ge1-\sqrt\eta,
\]
we obtain, since $\beta\le3+1/r\le4$,
\[
\begin{aligned}
  &\sqrt{\frac{1}{2(1+\eta)}
  \left(\binom q2-\frac{5qr}{4}\right)}\\
  &\qquad\ge
  r(1-\sqrt\eta)\sqrt{\frac{2\beta^2-5\beta}{8}}
  -\sqrt r.
\end{aligned}
\]
Moreover, $5-3\beta/2\ge0$ for all sufficiently large $r$.  Hence
\[
  m\ge r(1-\sqrt\eta)F(\beta)-5\sqrt r-5.
\]

A direct differentiation shows that $F$ is increasing on $(5/2,\infty)$:
after moving the negative term in $F'(\beta)$ to the other side and
squaring, the claim is equivalent to
$256\beta^2-640\beta+625>0$, whose leading coefficient is positive and
whose discriminant is negative.  Also,
$3\beta_0^2-8\beta_0-1=0$ and $\beta_0>1$, so
\[
  F(\beta)\ge F(\beta_0)
  =\frac{15}{4}-\frac{\beta_0}{4}=c.
\]
Consequently,
\[
  m\ge (1-\sqrt\eta)cr-5\sqrt r-5.
\]
Our choice of $\eta$ gives $c\sqrt\eta\le\varepsilon/2$, and for all
sufficiently large $r$ we have $5\sqrt r+5\le\varepsilon r/2$.
Thus $m\ge(c-\varepsilon)r$, as required.
\end{proof}

\end{document}
